% ============================================================
%  Guía de aplicación: del dato a la recomendación
%  Semana 12 · Clase 24: MasterClass (Recapitulación del Módulo 4)
%  Curso: Análisis y Visualización de Datos para la Toma de Decisiones
%  Datos: datos/colibri_envios.csv (simulados; empresa ficticia), generados
%  con ../../../scripts/semana12-clase24_datos.py
%  Las cifras de control se verificaron en Python (scikit-learn) y en
%  Orange 3.41: coinciden.
%  Compilar con: xelatex guia-aplicacion-modulo4.tex (dos veces)
% ============================================================
\documentclass[11pt]{article}

\usepackage[a4paper,margin=2cm]{geometry}
\usepackage{fontspec}
\usepackage{amsmath}
\usepackage{amssymb}
\usepackage[spanish,es-noshorthands,es-tabla]{babel}
\usepackage{xcolor}
\usepackage{graphicx}
\usepackage{booktabs}
\usepackage{array}
\usepackage{tabularx}
\usepackage{enumitem}
\usepackage{fancyvrb}
\usepackage{needspace}
\usepackage{tcolorbox}
\usepackage{tikz}
\usepackage{fancyhdr}
\usepackage{hyperref}

\tcbuselibrary{skins,breakable}
\usetikzlibrary{arrows.meta,positioning,calc,babel}

% ---------------- Tipografía ----------------
\IfFontExistsTF{Aptos}{\setmainfont{Aptos}}{\setmainfont{Lato}}
\setmonofont{DejaVu Sans Mono}[Scale=0.82]

% ---------------- Paleta Esumer ----------------
\definecolor{esumerNavy}{HTML}{262E66}
\definecolor{esumerNavyDark}{HTML}{1E4093}
\definecolor{esumerTeal}{HTML}{0CB3B6}
\definecolor{esumerSky}{HTML}{26A1DB}
\definecolor{esumerOrange}{HTML}{EF7C2C}
\definecolor{esumerGreen}{HTML}{94C122}
\definecolor{esumerMagenta}{HTML}{AF348B}
\definecolor{esumerGray}{HTML}{E8E8E8}
\definecolor{esumerBg}{HTML}{F7F8FA}

\hypersetup{colorlinks=true, linkcolor=esumerNavyDark, urlcolor=esumerNavyDark}
\graphicspath{{../../../plantilla/assets/}}

\newcommand{\urldatos}{https://github.com/craljimenez/An-lisis_y_Visualizacion_de_Datos_para_la_Toma_de_Decisiones/raw/main/clases/semana12-clase24-masterclass-recapitulacion-modulo4/datos/colibri_envios.csv}
\newcommand{\urlorange}{https://orangedatamining.com/download/}
\newcommand{\urlcolab}{https://colab.research.google.com/}
\newcommand{\urlformula}{https://orangedatamining.com/widget-catalog/transform/formula/}

% ---------------- Encabezado / pie de página ----------------
\pagestyle{fancy}
\fancyhf{}
\renewcommand{\headrulewidth}{0pt}
\fancyfoot[L]{\small\color{esumerNavyDark} Guía de aplicación del Módulo 4 \textperiodcentered\ Clase 24}
\fancyfoot[R]{\small\color{esumerNavyDark} Cristian A. Jiménez C. \quad\textperiodcentered\quad \thepage}
\setlength{\parindent}{0pt}
\setlength{\parskip}{4pt}
\setlength{\emergencystretch}{1.5em}

% ---------------- Código ----------------
\fvset{fontsize=\small}
\newtcolorbox{codigo}[1][]{
  colback=esumerBg, colframe=esumerNavy!35, boxrule=0.5pt, arc=1mm,
  left=6pt, right=6pt, top=2pt, bottom=2pt, fonttitle=\bfseries\small,
  coltitle=esumerNavy, colbacktitle=esumerGray, #1
}

% ---------------- Cajas reutilizables ----------------
\newtcolorbox{recuerda}[1][Lo viste en clase]{
  colback=esumerBg, colframe=esumerTeal, boxrule=0.6pt, arc=1mm,
  fonttitle=\bfseries\small, coltitle=white, colbacktitle=esumerTeal,
  title=#1, left=6pt, right=6pt, top=3pt, bottom=3pt
}
\newtcolorbox{cuidado}[1][Cuidado]{
  colback=white, colframe=esumerOrange, boxrule=0.6pt, arc=1mm,
  fonttitle=\bfseries\small, coltitle=white, colbacktitle=esumerOrange,
  title=#1, left=6pt, right=6pt, top=3pt, bottom=3pt
}
\newtcolorbox{control}[1][Cifras de control: compara antes de seguir]{
  colback=white, colframe=esumerGreen!80!black, boxrule=0.6pt, arc=1mm,
  fonttitle=\bfseries\small, coltitle=white, colbacktitle=esumerGreen!80!black,
  title=#1, left=6pt, right=6pt, top=3pt, bottom=3pt
}
\newtcolorbox{orange}[1][En Orange]{
  colback=white, colframe=esumerSky, boxrule=0.6pt, arc=1mm,
  fonttitle=\bfseries\small, coltitle=white, colbacktitle=esumerSky,
  title=#1, left=6pt, right=6pt, top=3pt, bottom=3pt
}

% Opción de Orange: nombre en español (nombre en inglés)
\newcommand{\menu}[2]{\textbf{#1} {\color{black!55}(#2)}}
\newcommand{\widget}[1]{\leavevmode{\color{esumerNavyDark}\textbf{#1}}}
\newcommand{\vez}[1]{\leavevmode\mbox{\color{esumerTeal}\textbf{[#1]}}}

% ---------------- Encabezado de cada etapa ----------------
\newcommand{\etapa}[3]{%
  \clearpage\noindent
  \begin{tikzpicture}[baseline=(n.base)]
    \node[circle, fill=esumerOrange, text=white, font=\bfseries\large,
          minimum size=0.85cm, inner sep=0pt] (n) {#1};
  \end{tikzpicture}\hspace{8pt}%
  {\Large\bfseries\color{esumerNavy} #2}\par
  \vspace{-2pt}{\color{esumerOrange}\rule{\textwidth}{1pt}}\par
  \vspace{2pt}
  {\color{esumerNavyDark}\textbf{Objetivo:}} \textit{#3}\par
  \vspace{2pt}
}
\newcommand{\subtitulo}[1]{\par\needspace{3\baselineskip}\vspace{4pt}{\large\bfseries\color{esumerNavyDark} #1}\par}

\newlist{checklist}{itemize}{1}
\setlist[checklist]{label=$\square$, itemsep=1pt, topsep=2pt, leftmargin=1.4em}

\tikzset{
  every node/.append style={execute at begin node={\hyphenpenalty=10000\exhyphenpenalty=10000\relax}},
  flecha/.style={-{Stealth[length=2.4mm]}, line width=0.9pt, color=esumerNavy},
  num/.style={circle, fill=esumerOrange, text=white, font=\bfseries\small,
              minimum size=0.6cm, inner sep=0pt},
}

% ============================================================
\begin{document}

\begin{minipage}[c]{0.18\textwidth}
  \includegraphics[height=1.15cm,width=\linewidth,keepaspectratio]{esumer-logo.png}
\end{minipage}%
\begin{minipage}[c]{0.80\textwidth}
  \raggedleft
  {\Large\bfseries\color{esumerNavy} Guía de aplicación: del dato a la recomendación}\\[3pt]
  {\small\color{esumerNavyDark} Semana 12 \textperiodcentered\ Clase 24 \textperiodcentered\ MasterClass del Módulo 4}
\end{minipage}

\vspace{4pt}
{\color{esumerOrange}\rule{\textwidth}{2pt}}
\vspace{6pt}

\textbf{Para qué es esta guía.} En el Módulo 4 aprendiste a pasar de
describir a predecir (Clase 19), a construir y evaluar una regresión lineal
(Clase 20), a contar una historia (Clase 21), a estimar el impacto de una
decisión (Clase 22) y a presentarla (Clase 23). Hoy haces todo el recorrido,
de principio a fin, con un caso nuevo. La guía te dice qué hacer y cómo
comprobar que vas bien; cuando dudes, las pistas como \vez{Clase 20} te
dicen dónde lo viste.

\begin{recuerda}[El caso: Colibrí Envíos S.A.S.]
\small Colibrí Envíos es una empresa de mensajería en el Valle de Aburrá:
recoge pedidos de tiendas y los entrega en moto o en bicicleta. Hoy, la app
le promete a cada cliente que su pedido llegará en \textbf{45 minutos},
sin importar la distancia ni la hora. Cuando llega tarde, la empresa le
regala al cliente un \textbf{bono de 8.000 pesos}. Tienes los 1500 pedidos
entregados de enero a junio de 2026.

\vspace{2pt}
\footnotesize\itshape La empresa es ficticia y los datos son simulados con
fines académicos.
\end{recuerda}

\textbf{Al terminar la sesión deberías tener:}
\begin{enumerate}[itemsep=1pt, topsep=2pt]
  \item La pregunta predictiva escrita con la plantilla de la Clase 19.
  \item La tabla de características: cuáles entran, cuáles no y por qué.
  \item Un modelo de regresión lineal evaluado con datos de prueba, en
    Python o en Orange.
  \item El impacto estimado de una nueva política de promesa, en pesos y
    con sus supuestos.
  \item El guion de una presentación de cinco minutos para la gerencia.
\end{enumerate}

\textbf{Agenda sugerida para las 3 horas}

\begin{tabular}{@{}>{\raggedright\arraybackslash}p{1.4cm}>{\raggedright\arraybackslash}p{6.9cm}>{\raggedright\arraybackslash}p{7.6cm}@{}}
\toprule
\textbf{Hora} & \textbf{Etapas} & \textbf{Producto al final de la hora}\\
\midrule
Hora 1 & 1. Pregunta \textperiodcentered\ 2. Características & Pregunta predictiva y tabla de características.\\
Hora 2 & 3. Modelo \textperiodcentered\ 4. Límites & Tres modelos comparados y sus límites.\\
Hora 3 & 5. Impacto \textperiodcentered\ 6. Comunicar & Impacto en pesos y guion de la presentación.\\
\bottomrule
\end{tabular}

\vspace{6pt}
\textbf{Elige tu ruta.} Las etapas 3 a 5 se pueden hacer en
\href{\urlcolab}{\color{esumerOrange}Python (Google Colab)} o en
\href{\urlorange}{\color{esumerOrange}Orange}. Las cifras de control son
\textbf{las mismas en las dos}. Se puede trabajar en parejas o de forma
individual. Descarga el archivo
\href{\urldatos}{\color{esumerOrange}\texttt{colibri\_envios.csv}}.

% ============================================================
\clearpage
\section*{\color{esumerNavy}La carta de la gerencia}

\begin{tcolorbox}[colback=white, colframe=esumerNavy, boxrule=0.8pt, arc=1mm,
  left=10pt, right=10pt, top=6pt, bottom=6pt]
\small
\textit{Hola. Cerca del 29\% de nuestros pedidos llega después de los 45
minutos que prometemos, y cada uno nos cuesta un bono de 8.000 pesos y un
cliente molesto. Para el comité de diciembre necesito saber:}

\begin{enumerate}[label=\textbf{P\arabic*.}, leftmargin=2.2em, itemsep=3pt, topsep=4pt]
  \item \textit{¿Podemos estimar, al momento de recibir un pedido, cuántos
    minutos va a tardar? ¿Con qué error?}
  \item \textit{¿Qué es lo que más hace tardar un pedido?}
  \item \textit{Si en vez de prometer 45 minutos a todos, la app prometiera
    un tiempo según cada pedido, ¿cuántos pedidos llegarían tarde? ¿Cuánto
    nos ahorraríamos en bonos?}
  \item \textit{¿Qué riesgos tiene usar un modelo para esto?}
\end{enumerate}

\vspace{4pt}
\textit{Gracias.}\\
\textbf{Gerencia de operaciones, Colibrí Envíos}
\end{tcolorbox}

\subtitulo{Los datos: \texttt{colibri\_envios.csv}}

Una fila por pedido entregado (1500 filas).

\small
\renewcommand{\arraystretch}{1.12}
\begin{tabularx}{\textwidth}{@{}>{\raggedright\arraybackslash}p{4.0cm} >{\raggedright\arraybackslash}X@{}}
\toprule
\textbf{Columna} & \textbf{Qué contiene}\\
\midrule
\texttt{id\_pedido} & Número del pedido.\\
\texttt{fecha}, \texttt{hora} & Día y hora en que se hizo el pedido (la hora, de 8 a 19).\\
\texttt{zona} & Zona de entrega: Centro, El Poblado, Laureles, Envigado o Bello.\\
\texttt{distancia\_km} & Distancia de la tienda al cliente, en kilómetros.\\
\texttt{num\_paquetes}, \texttt{peso\_kg} & Cuántos paquetes lleva el pedido y cuánto pesan en total.\\
\texttt{vehiculo} & Moto o Bicicleta.\\
\texttt{lluvia} & 1 si llovía cuando se hizo el pedido; 0 si no.\\
\texttt{cliente\_nuevo} & 1 si es el primer pedido del cliente; 0 si no.\\
\texttt{minutos\_entrega} & \textbf{Minutos desde el pedido hasta la entrega.}\\
\texttt{entrega\_tarde} & 1 si el pedido tardó más de 45 minutos.\\
\texttt{calificacion} & Calificación que dio el cliente al recibir, de 1 a 5.\\
\bottomrule
\end{tabularx}\par
\renewcommand{\arraystretch}{1.0}
\normalsize

\subtitulo{El recorrido completo}

\begin{center}
\begin{tikzpicture}[
  et/.style={draw=esumerNavy, line width=0.8pt, rounded corners=3pt,
             fill=white, minimum width=5.0cm, minimum height=2.3cm,
             text width=4.7cm, align=center, inner sep=3pt, font=\scriptsize},
  clase/.style={fill=esumerTeal, text=white, font=\scriptsize\bfseries,
                rounded corners=2pt, inner sep=2.5pt},
]
  \node[et] (e1) at (2.85,0) {{\bfseries\small La pregunta predictiva}\\[3pt]Decisión, objetivo, fila, tipo de tarea};
  \node[et] (e2) at (8.5,0) {{\bfseries\small Características}\\[3pt]Extraer, medir la relación, detectar la fuga};
  \node[et] (e3) at (14.15,0) {{\bfseries\small El modelo}\\[3pt]Entrenamiento y prueba, línea base, regresión simple y múltiple};
  \node[et] (e4) at (14.15,-3.1) {{\bfseries\small Los límites}\\[3pt]Errores grandes, extrapolación, una prueba de fuga};
  \node[et] (e5) at (8.5,-3.1) {{\bfseries\small El impacto}\\[3pt]Una nueva promesa: pedidos tarde y bonos ahorrados};
  \node[et] (e6) at (2.85,-3.1) {{\bfseries\small Comunicar}\\[3pt]El mensaje, el guion y el pitch de cinco minutos};
  \foreach \i in {1,...,6} {\node[num] at (e\i.north west) {\i};}
  \node[clase, anchor=south east] at ($(e1.north east)+(-0.15,-0.02)$) {Clase 19};
  \node[clase, anchor=south east] at ($(e2.north east)+(-0.15,-0.02)$) {Clase 19};
  \node[clase, anchor=south east] at ($(e3.north east)+(-0.15,-0.02)$) {Clase 20};
  \node[clase, anchor=south east] at ($(e4.north east)+(-0.15,-0.02)$) {Clase 20};
  \node[clase, anchor=south east] at ($(e5.north east)+(-0.15,-0.02)$) {Clase 22};
  \node[clase, anchor=south east] at ($(e6.north east)+(-0.15,-0.02)$) {Clases 21 y 23};
  \draw[flecha] (e1) -- (e2);
  \draw[flecha] (e2) -- (e3);
  \draw[flecha] (e3) -- (e4);
  \draw[flecha] (e4) -- (e5);
  \draw[flecha] (e5) -- (e6);
\end{tikzpicture}
\end{center}

% ============================================================
\etapa{1}{La pregunta predictiva}{convertir la carta de la gerencia en una pregunta que un modelo pueda responder.}

Completa la plantilla de la Clase 19 con el caso de Colibrí Envíos.
\vez{Clase 19}

\renewcommand{\arraystretch}{2.0}
\begin{tabularx}{\textwidth}{@{}>{\raggedright\arraybackslash}p{4.2cm} >{\raggedright\arraybackslash}X@{}}
\toprule
\textbf{Para decidir} (qué haremos) & \\ \hline
\textbf{quiero predecir} (el objetivo) & \\ \hline
\textbf{de cada} (la fila) & \\ \hline
\textbf{usando} (lo que sé en ese momento) & \\ \hline
\textbf{con anticipación de} & \\
\bottomrule
\end{tabularx}\par
\renewcommand{\arraystretch}{1.0}

\vspace{4pt}
Responde también:

\begin{enumerate}[itemsep=2pt, topsep=2pt]
  \item ¿Es un problema de clasificación o de regresión? ¿Por qué?
  \item ¿Cuál sería la línea base: qué hace hoy la empresa, sin modelo?
  \item ¿En qué momento exacto se usaría la predicción? Esto decide qué
    columnas puedes usar en la etapa 2.
\end{enumerate}

\begin{control}[Para comparar con tu respuesta]
\small \textbf{Para decidir} qué tiempo de entrega prometer en la app,
\textbf{quiero predecir} los minutos que tardará \textbf{cada pedido},
\textbf{usando} lo que se sabe al recibirlo (distancia, paquetes,
vehículo, hora y si llueve), \textbf{con anticipación de} todo el tiempo
de la entrega: se predice en el momento del pedido.

Es \textbf{regresión}, porque el objetivo (\texttt{minutos\_entrega}) es un
número. La línea base es la promesa de hoy: 45 minutos para todos o, para
medir el error del modelo, el promedio de minutos del entrenamiento.
\end{control}

\begin{cuidado}[El momento de la predicción manda]
\small Si la app debe mostrar la promesa cuando el cliente hace el pedido,
solo vale lo que se sabe en ese instante. Esa frase te va a sacar de apuros
en la etapa 2. \vez{Clase 19, fuga de información}
\end{cuidado}

% ============================================================
\etapa{2}{Características: extraer y seleccionar}{decidir con qué columnas aprende el modelo, y crear las que hagan falta.}

\subtitulo{Extraer dos columnas nuevas}

\begin{enumerate}[itemsep=2pt, topsep=2pt]
  \item \textbf{\texttt{hora\_pico}}: 1 si el pedido se hizo a las 12, 13,
    18 o 19 horas (almuerzo y salida del trabajo); 0 si no. La hora, tal
    cual, es un número del 8 al 19, y el tráfico no sube ni baja en línea
    recta con ella.
  \item \textbf{\texttt{es\_bicicleta}}: 1 si el vehículo es Bicicleta; 0 si
    es Moto. El modelo lineal necesita números.
\end{enumerate}

\begin{codigo}[title=En Python (Colab)]
\begin{Verbatim}
import pandas as pd

pedidos = pd.read_csv("colibri_envios.csv")
pedidos["hora_pico"] = pedidos["hora"].isin([12, 13, 18, 19]).astype(int)
pedidos["es_bicicleta"] = (pedidos["vehiculo"] == "Bicicleta").astype(int)
\end{Verbatim}
\end{codigo}

\begin{orange}
\small \widget{File} con el CSV. En la lista de columnas, cambia el tipo de
\texttt{lluvia}, \texttt{cliente\_nuevo} y \texttt{entrega\_tarde} a
\menu{numérico}{numeric} (Clase 20, paso 1) y pon \texttt{id\_pedido} en
\menu{omitir}{skip}. Luego un widget \widget{Formula}
(\href{\urlformula}{\color{esumerOrange}documentación}): \menu{Nuevo}{New}
$>$ \menu{Numérico}{Numeric}, con el nombre y la expresión:

\vspace{2pt}
\texttt{hora\_pico}: \texttt{1 if hora in (12, 13, 18, 19) else 0}\\
\texttt{es\_bicicleta}: \texttt{1 if vehiculo == 'Bicicleta' else 0}

\vspace{2pt}
Pulsa \menu{Enviar}{Send}.
\end{orange}

\subtitulo{Medir la relación con el objetivo}

\begin{codigo}[title=En Python]
\begin{Verbatim}
columnas = ["distancia_km", "num_paquetes", "peso_kg", "lluvia", "cliente_nuevo",
            "hora", "hora_pico", "es_bicicleta", "calificacion", "entrega_tarde"]
pedidos[columnas].corrwith(pedidos["minutos_entrega"]).round(3)
\end{Verbatim}
\end{codigo}

\begin{orange}
\small \widget{Select Columns} con \texttt{minutos\_entrega} en
\menu{Objetivo}{Target} y las demás columnas numéricas en
\menu{Variables}{Features}. Luego \widget{Rank}: con un objetivo numérico
muestra \textbf{Univar. reg.}, que ordena las columnas igual que la
correlación.
\end{orange}

\begin{control}
\small Correlación con \texttt{minutos\_entrega}:
\vspace{2pt}

\renewcommand{\arraystretch}{1.1}
\begin{tabularx}{\linewidth}{@{}>{\raggedright\arraybackslash}X >{\raggedleft\arraybackslash}p{1.6cm} @{\hspace{1.0cm}} >{\raggedright\arraybackslash}X >{\raggedleft\arraybackslash}p{1.6cm}@{}}
\toprule
\texttt{entrega\_tarde} & 0,770 & \texttt{num\_paquetes} & 0,194\\
\texttt{calificacion} & $-0{,}767$ & \texttt{peso\_kg} & 0,177\\
\texttt{distancia\_km} & 0,652 & \texttt{hora} & 0,141\\
\texttt{hora\_pico} & 0,347 & \texttt{es\_bicicleta} & 0,132\\
\texttt{lluvia} & 0,212 & \texttt{cliente\_nuevo} & $-0{,}013$\\
\bottomrule
\end{tabularx}\par
\vspace{2pt}
\texttt{hora\_pico} se relaciona con el tiempo más del doble que
\texttt{hora} (0,347 frente a 0,141): extraerla valió la pena.
\end{control}

\subtitulo{Decidir: la tabla de características}

Pasa cada columna por las cinco preguntas de la Clase 19 y completa la
tabla. La primera fila es un ejemplo. \vez{Clase 19}

\small
\renewcommand{\arraystretch}{1.55}
\begin{tabularx}{\textwidth}{@{}>{\raggedright\arraybackslash}p{2.9cm} >{\centering\arraybackslash}p{1.6cm} >{\centering\arraybackslash}p{1.6cm} >{\raggedright\arraybackslash}p{2.0cm} >{\raggedright\arraybackslash}X@{}}
\toprule
\textbf{Columna} & \textbf{¿La tendré?} & \textbf{¿Aporta?} & \textbf{Decisión} & \textbf{Por qué}\\
\midrule
{\itshape distancia\_km} & {\itshape Sí} & {\itshape Sí} & {\itshape Entra} & {\itshape Se conoce al pedir y es la que más se relaciona.}\\ \hline
num\_paquetes & & & & \\ \hline
peso\_kg & & & & \\ \hline
lluvia & & & & \\ \hline
cliente\_nuevo & & & & \\ \hline
hora / hora\_pico & & & & \\ \hline
es\_bicicleta & & & & \\ \hline
calificacion & & & & \\ \hline
entrega\_tarde & & & & \\ \hline
id\_pedido & & & & \\
\bottomrule
\end{tabularx}\par
\renewcommand{\arraystretch}{1.0}
\normalsize

\begin{control}[Para comparar con tu tabla]
\small\raggedright
\textbf{Entran}: distancia, paquetes, lluvia, \texttt{hora\_pico} y
\texttt{es\_bicicleta}.

\textbf{Fuera}:
\begin{itemize}[itemsep=0pt, topsep=1pt]
  \item \texttt{calificacion} y \texttt{entrega\_tarde}: aunque encabezan
    la tabla, se conocen \textbf{después} de entregar (fuga de información).
  \item \texttt{peso\_kg}: repite lo que dice el número de paquetes
    (correlación entre ellas: 0,95).
  \item \texttt{hora}: ya la reemplaza \texttt{hora\_pico}.
  \item \texttt{cliente\_nuevo}: no se relaciona con el tiempo.
  \item \texttt{id\_pedido}: es un identificador.
\end{itemize}
\end{control}

% ============================================================
\etapa{3}{El modelo}{entrenar una regresión lineal y saber cuánto mejora a no hacer nada.}

\subtitulo{Separar, línea base y dos modelos}

\begin{codigo}[title=En Python]
\begin{Verbatim}
from sklearn.model_selection import train_test_split
from sklearn.linear_model import LinearRegression
from sklearn.metrics import mean_absolute_error, r2_score

entrenamiento, prueba = train_test_split(pedidos, test_size=0.2, random_state=42)
y_train = entrenamiento["minutos_entrega"]
y_test = prueba["minutos_entrega"]

base = [y_train.mean()] * len(prueba)
print("Línea base, MAE:", round(mean_absolute_error(y_test, base), 2))

columnas = ["distancia_km", "num_paquetes", "lluvia", "hora_pico", "es_bicicleta"]
multiple = LinearRegression()
multiple.fit(entrenamiento[columnas], y_train)
pred = multiple.predict(prueba[columnas])
print("Múltiple, MAE:", round(mean_absolute_error(y_test, pred), 2),
      "R2:", round(r2_score(y_test, pred), 3))
\end{Verbatim}
\end{codigo}

\begin{orange}
\small El mismo flujo de la Clase 20: \widget{Data Sampler} al 20\%, con
\menu{Muestreo replicable}{Replicable sampling};
\widget{Constant} y \widget{Linear Regression} conectados a
\widget{Test and Score}, con \menu{Datos restantes}{Remaining Data} como
entrenamiento y \menu{Muestra de datos}{Data Sample} como
\menu{Datos de prueba}{Test Data}, y \menu{Probar con datos de
prueba}{Test on test data}. Los modelos se cambian en
\widget{Select Columns}. \vez{Clase 20, pasos 3 a 7}
\end{orange}

Prueba tres modelos y compáralos con la línea base: solo
\texttt{distancia\_km}; las cinco columnas, pero con \texttt{hora} en vez de
\texttt{hora\_pico}; y las cinco columnas elegidas.

\begin{control}
\small 1200 pedidos de entrenamiento y 300 de prueba. Con los 300 de prueba:
\vspace{2pt}

\renewcommand{\arraystretch}{1.15}
\begin{tabularx}{\linewidth}{@{}>{\raggedright\arraybackslash}X >{\raggedleft\arraybackslash}p{2.0cm} >{\raggedleft\arraybackslash}p{2.0cm} >{\raggedleft\arraybackslash}p{2.0cm}@{}}
\toprule
\textbf{Modelo} & \textbf{MAE} & \textbf{RMSE} & \textbf{$R^2$}\\
\midrule
Línea base (promedio: 39,9 minutos) & 10,40 & 12,95 & $-0{,}004$\\
Simple: solo la distancia & 7,56 & 9,90 & 0,413\\
Cinco columnas, con \texttt{hora} & 6,18 & 8,66 & 0,551\\
Cinco columnas, con \texttt{hora\_pico} & \textbf{5,02} & \textbf{7,81} & \textbf{0,635}\\
\bottomrule
\end{tabularx}\par
\vspace{2pt}
Extraer \texttt{hora\_pico} bajó el error de 6,18 a 5,02 minutos. El
modelo final se equivoca en unos 5 minutos por pedido; la línea base, en
unos 10.
\end{control}

\subtitulo{Los coeficientes, en lenguaje de negocio}

\begin{control}[Coeficientes del modelo final]
\small
\renewcommand{\arraystretch}{1.15}
\begin{tabularx}{\linewidth}{@{}>{\raggedright\arraybackslash}p{2.8cm} >{\raggedleft\arraybackslash}p{1.8cm} @{\hspace{0.6cm}} >{\raggedright\arraybackslash}X@{}}
\toprule
\textbf{Variable} & \textbf{Coeficiente} & \textbf{Con lo demás igual\ldots}\\
\midrule
intercepto & 13,46 & \\
\texttt{distancia\_km} & 2,99 & cada kilómetro suma unos 3 minutos.\\
\texttt{num\_paquetes} & 1,86 & cada paquete más, unos 2 minutos.\\
\texttt{lluvia} & 6,70 & si llueve, unos 7 minutos más.\\
\texttt{hora\_pico} & 9,14 & en hora pico, unos 9 minutos más.\\
\texttt{es\_bicicleta} & 4,38 & en bicicleta, unos 4 minutos más que en moto.\\
\bottomrule
\end{tabularx}\par
\end{control}

Con esto respondes P1 y P2 de la carta. ¿Qué es lo que más pesa? Ojo: un
coeficiente grande no es más importante por sí solo, depende de las
unidades. Un pedido típico recorre 4,8 km: la distancia aporta unos 14
minutos, más que la hora pico. \vez{Clase 20}

% ============================================================
\etapa{4}{Los límites del modelo}{saber cuándo no usarlo antes de que alguien te lo pregunte.}

\begin{enumerate}[itemsep=3pt, topsep=2pt]
  \item \textbf{Los errores más grandes.} Ordena los pedidos de prueba por
    su residuo (real menos predicho). \vez{Clase 20, paso 9}
\begin{codigo}
\begin{Verbatim}
prueba = prueba.copy()
prueba["prediccion"] = pred.round(0)
prueba["residuo"] = prueba["minutos_entrega"] - prueba["prediccion"]
prueba.sort_values("residuo", ascending=False).head(4)
\end{Verbatim}
\end{codigo}
    En Orange: \widget{Predictions} con la muestra de prueba y el modelo,
    ordenando por la diferencia (recuerda que allí el signo va al revés).
  \item \textbf{Extrapolar.} ¿Qué predice el modelo para un pedido de 25 km,
    con un paquete, sin lluvia, fuera de hora pico y en moto? ¿Entre qué
    distancias aprendió?
  \item \textbf{La prueba de la fuga.} Agrega \texttt{calificacion} a las
    columnas del modelo y mira qué pasa con el $R^2$.
\end{enumerate}

\begin{control}
\small \textbf{Errores grandes.} El mayor es el pedido 11235: el modelo
predijo unos 38 minutos y tardó 83 (residuo de 45). Le siguen pedidos de 65
minutos predichos en 25 y 30. Son pedidos cortos que tardaron mucho más de
lo normal: algo pasó (una dirección errada, un cliente que no contestaba)
que no está en ninguna columna. \textbf{El modelo no ve las novedades.}

\textbf{Extrapolar.} Predice unos 90 minutos para 25 km, pero aprendió con
pedidos de 0,5 a 18 km: fuera de ese rango, no se usa.

\textbf{La fuga.} Con \texttt{calificacion}, el $R^2$ sube de 0,635 a 0,759
y el MAE baja a 4,48. Parece mejor, pero la calificación se da después de
entregar: el modelo no se podría usar al recibir el pedido. Un resultado
demasiado bueno casi siempre es una fuga. \vez{Clase 19}
\end{control}

\begin{recuerda}[Para la presentación]
\small Estos tres hallazgos son tu diapositiva de limitaciones (Clase 23):
``El modelo sirve para pedidos normales de hasta 18 km; no anticipa
novedades como una dirección errada''. Responde P4.
\end{recuerda}

% ============================================================
\etapa{5}{El impacto}{estimar cuánto cambia la decisión con el modelo, en pedidos tarde y en pesos.}

\textbf{La nueva política.} En vez de prometer 45 minutos a todos, la app
promete \textbf{la predicción del modelo más 10 minutos} de margen (los 10
minutos cubren el error típico de unos 5, con holgura). ¿Cuántos pedidos de
prueba habrían llegado tarde con cada política?

\begin{codigo}[title=En Python]
\begin{Verbatim}
real = prueba["minutos_entrega"]
promesa_modelo = pred + 10          # la predicción sin redondear, más 10 minutos

print("Tarde con 45 minutos:", (real > 45).sum())
print("Tarde con el modelo:", (real > promesa_modelo).sum())
print("Promesa promedio con el modelo:", round(promesa_modelo.mean(), 1))
print("Promesas de más de 60 minutos:", (promesa_modelo > 60).sum())
\end{Verbatim}
\end{codigo}

\begin{orange}
\small Conecta un \widget{Formula} a la salida de \widget{Predictions} (la
de los datos de prueba). Crea dos columnas numéricas:
\texttt{tarde\_hoy} con \texttt{1 if minutos\_entrega > 45 else 0} y
\texttt{tarde\_modelo} con
\texttt{1 if minutos\_entrega > Linear\_Regression + 10 else 0}. En la
expresión, el espacio de ``Linear Regression'' se escribe con guion bajo.
Luego el widget \widget{Column Statistics}: el promedio de cada columna es
la proporción de pedidos tarde.
\end{orange}

\begin{control}
\small Con los 300 pedidos de prueba:
\vspace{2pt}

\renewcommand{\arraystretch}{1.15}
\begin{tabularx}{\linewidth}{@{}>{\raggedright\arraybackslash}X >{\raggedleft\arraybackslash}p{2.8cm} >{\raggedleft\arraybackslash}p{2.8cm}@{}}
\toprule
& \textbf{Hoy: 45 minutos} & \textbf{Modelo + 10}\\
\midrule
Pedidos tarde & 85 (28,3\%) & 14 (4,7\%)\\
Promesa promedio & 45,0 minutos & 48,9 minutos\\
Promesas de más de 60 minutos & 0 & 50\\
\bottomrule
\end{tabularx}\par
\end{control}

\subtitulo{De pedidos a pesos}

Colibrí Envíos entrega unos \textbf{250 pedidos al mes} (1500 en seis
meses). Con los supuestos escritos, estima el ahorro en bonos:

\renewcommand{\arraystretch}{1.25}
\begin{tabularx}{\textwidth}{@{}>{\raggedright\arraybackslash}X >{\raggedleft\arraybackslash}p{3.0cm} >{\raggedleft\arraybackslash}p{3.0cm}@{}}
\toprule
& \textbf{Hoy} & \textbf{Con el modelo}\\
\midrule
Pedidos tarde al mes ($250 \times$ la proporción de la prueba) & & \\
Bonos al mes ($\times$ 8.000 pesos) & & \\
\midrule
\textbf{Ahorro mensual} & \multicolumn{2}{c}{} \\
\bottomrule
\end{tabularx}\par
\renewcommand{\arraystretch}{1.0}

\begin{control}[Para comparar]
\small Hoy: $250 \times 85/300 \approx 71$ pedidos tarde y unos
\textbf{567.000} pesos en bonos al mes. Con el modelo: $250 \times
14/300 \approx 12$ pedidos y unos \textbf{93.000} pesos. Ahorro: unos
\textbf{470.000 pesos al mes}.

\textbf{Supuestos:} el volumen se mantiene en 250 pedidos; la proporción
de pedidos tarde de la prueba vale para los meses que vienen; el bono sigue
en 8.000 pesos. \textbf{Lo que falta saber:} si una promesa más larga (en
promedio 3,9 minutos más, y 50 de cada 300 por encima de una hora) hace que
algunos clientes no pidan. Ese costo no está en los datos. \vez{Clase 22}
\end{control}

\subtitulo{La matriz indicador, hallazgo, decisión, impacto}

Resume las etapas 3 a 5 en la matriz de la Clase 22, con dos o tres filas.
Por ejemplo: indicador, \emph{porcentaje de pedidos tarde}; hallazgo,
\emph{28,3\% con la promesa fija}; decisión, \emph{prometer según el modelo
más 10 minutos}; impacto, \emph{unos 470.000 pesos al mes en bonos}.
\vez{Clase 22}

% ============================================================
\etapa{6}{Comunicar}{contarle a la gerencia, en cinco minutos, qué encontraste y qué propones.}

\subtitulo{El mensaje en una frase}

Escríbela con la fórmula de la Clase 21: qué pasó, por qué importa y qué
propones. \vez{Clase 21}

\renewcommand{\arraystretch}{3.0}
\begin{tabularx}{\textwidth}{@{}>{\raggedright\arraybackslash}X@{}}
\toprule
\\
\bottomrule
\end{tabularx}\par
\renewcommand{\arraystretch}{1.0}

\subtitulo{El guion: cinco diapositivas}

Escribe solo los títulos, que concluyan. Leídos en orden, deben contar la
historia completa.\par\hfill\vez{Clases 21 y 23}

\renewcommand{\arraystretch}{2.2}
\begin{tabularx}{\textwidth}{@{}>{\centering\arraybackslash}p{0.6cm} >{\raggedright\arraybackslash}p{3.4cm} >{\raggedright\arraybackslash}X@{}}
\toprule
& \textbf{Parte del pitch} & \textbf{Título de la diapositiva}\\
\midrule
1 & Apertura (30 s) & \\ \hline
2 & Evidencia y modelo (2 min) & \\ \hline
3 & Limitaciones & \\ \hline
4 & Recomendación e impacto (2 min) & \\ \hline
5 & Cierre: qué pedimos hoy (30 s) & \\
\bottomrule
\end{tabularx}\par
\renewcommand{\arraystretch}{1.0}

\begin{control}[Un ejemplo de títulos]
\small (1) Hoy, casi 3 de cada 10 pedidos llegan después de la hora
prometida. (2) Con la distancia, la hora y el clima, estimamos el tiempo
con un error de unos 5 minutos. (3) El modelo no anticipa novedades ni
sirve para más de 18 km. (4) Prometer según el modelo bajaría los pedidos
tarde a menos de 1 de cada 20 y ahorraría unos 470.000 pesos al mes.
(5) Pedimos aprobar un piloto de un mes en una zona.
\end{control}

\subtitulo{El ensayo}

En parejas, cuenta el pitch en cinco minutos; tu compañero hace de gerente
y te pregunta. Si trabajas solo, grábate. Revisa con la rúbrica de la
Clase 23: apertura, mensaje, modelo sin jerga, incertidumbre,
recomendación con impacto, tiempo y preguntas. \vez{Clase 23}

% ============================================================
\clearpage
\section*{\color{esumerNavy}Autoevaluación: ¿cubriste todo el módulo?}

\small
\renewcommand{\arraystretch}{1.3}
\begin{tabularx}{\textwidth}{@{}>{\raggedright\arraybackslash}p{2.7cm} >{\raggedright\arraybackslash}X >{\raggedright\arraybackslash}X >{\centering\arraybackslash}p{0.9cm}@{}}
\toprule
\textbf{Criterio} & \textbf{En proceso} & \textbf{Logrado} & \textbf{Yo}\\
\midrule
\textbf{Pregunta} & Falta la decisión o el momento de la predicción. & Decisión, objetivo, fila, datos disponibles y anticipación. & \\
\textbf{Características} & Usa todas las columnas, o se le coló una fuga. & Extrae \texttt{hora\_pico}, descarta la fuga, lo redundante y el identificador, y explica por qué. & \\
\textbf{Modelo} & Mide con los datos de entrenamiento o sin línea base. & Mide con la prueba, frente a la línea base, y lee los coeficientes en lenguaje de negocio. & \\
\textbf{Límites} & No los menciona. & Errores grandes, rango de uso y fuga, explicados. & \\
\textbf{Impacto} & Sin cifra o sin supuestos. & Pedidos tarde y pesos, con supuestos y con lo que falta saber. & \\
\textbf{Comunicar} & Describe el proceso, sin recomendación. & Mensaje en una frase, títulos que concluyen y una decisión pedida. & \\
\bottomrule
\end{tabularx}\par
\renewcommand{\arraystretch}{1.0}
\normalsize

\subtitulo{Si algo no cuadra}

\small
\renewcommand{\arraystretch}{1.2}
\begin{tabularx}{\textwidth}{@{}>{\raggedright\arraybackslash}p{5.8cm} >{\raggedright\arraybackslash}X@{}}
\toprule
\textbf{Lo que ves} & \textbf{Causa más probable}\\
\midrule
Todas las cifras del modelo son distintas. & Otro reparto: revisa \texttt{test\_size=0.2} y \texttt{random\_state=42}, o en Orange, 20\% con muestreo replicable.\\
El $R^2$ es de 0,759 o más. & Se coló \texttt{calificacion} o \texttt{entrega\_tarde} en las columnas.\\
El MAE es 6,18 y no 5,02. & Usaste \texttt{hora} en vez de \texttt{hora\_pico}.\\
En Orange, dos coeficientes por columna con signo contrario. & Faltó cambiar \texttt{lluvia} a numérico en \widget{File} (Clase 20).\\
\texttt{hora\_pico} es siempre 0 en Orange. & La expresión de \widget{Formula}: revisa los paréntesis y las comas de \texttt{(12, 13, 18, 19)}.\\
La fórmula de \texttt{tarde\_modelo} da error en Orange. & El nombre de la columna de predicción se escribe \texttt{Linear\_Regression}, con guion bajo.\\
\bottomrule
\end{tabularx}\par
\renewcommand{\arraystretch}{1.0}
\normalsize

\end{document}
